\documentclass[12pt]{article}

\usepackage{geometry}
\geometry{margin=1in}
\usepackage{titlesec}
\usepackage{enumitem}
\usepackage{booktabs}
\usepackage{array}
\usepackage{parskip}
\usepackage{hyperref}
\usepackage{xcolor}

\titleformat{\section}{\large\bfseries}{\thesection.}{0.5em}{}
\titleformat{\subsection}{\normalsize\bfseries}{\thesubsection}{0.5em}{}
\titleformat{\subsubsection}{\normalsize\itshape}{\thesubsubsection}{0.5em}{}

\newcommand{\experiment}[1]{\textcolor{blue}{\textit{Experiment: #1}}}
\newcommand{\contrib}[1]{\textbf{[#1]}}

\begin{document}

% ===========================================================================
% Title
% ===========================================================================
\begin{center}
    {\LARGE\bfseries Thesis Outline}\\[0.5em]
    {\large Reinforcement Learning for Fairness-Aware Synthetic Data Generation\\
    under Positive-Class Scarcity}\\[0.5em]
    {\normalsize Ethan Pigou\\
    Department of Electrical and Computer Engineering, Western University}
\end{center}

\vspace{1em}
\hrule
\vspace{1em}

% ===========================================================================
% Thesis Contributions
% ===========================================================================
\section*{Thesis Contributions}

\begin{description}[leftmargin=2.5cm, labelwidth=2cm]

    \item[\textbf{C1 --- FORGE method.}]
    A reinforcement learning agent generates synthetic minority-positive
    training samples in a PCA-compressed feature space, guided by a
    worst-group validation loss reward, improving group fairness under severe
    positive-class scarcity without modifying the downstream classifier
    objective. Demonstrated empirically across two primary datasets against six
    baseline methods in two feature-space settings.

    \item[\textbf{C2 --- Failure mode characterisation.}]
    Under positive-class scarcity, reweighting methods plateau due to a
    structural data constraint, and naive generative methods degrade fairness
    because they synthesise without downstream feedback. Both failure modes are
    characterised empirically through baseline comparison tables and spatially
    through PCA-space visualisation of synthetic point placement.

    \item[\textbf{C3 --- Governing design choices.}]
    Reward sharpness (sigmoid $k$) and classifier training intensity (epochs
    per episode) are the dominant design parameters for FORGE, and their
    optimal values are predictable from dataset-level reward signal diagnostics
    --- specifically episode return variance and deadzone fraction ---
    measurable during training. Demonstrated through a systematic grid search
    of 324 census and 243 Capture-24 configurations.

\end{description}

\vspace{0.5em}
\noindent\textbf{Relationship to prior work.}
Pigou et al.\ (IEEE CASSE 2025) demonstrated that RL-guided synthetic
generation outperforms GANs for class-level imbalance in human activity
recognition. This thesis extends the feedback-driven generation approach to
the group fairness setting under positive-class scarcity, replacing the
decision-boundary proximity reward with worst-group validation loss and
operating in a PCA-compressed feature space. Formal ablation studies and
multi-dataset evaluation are added.

\vspace{1em}
\hrule
\vspace{1em}

% ===========================================================================
% Chapter 1
% ===========================================================================
\section{Introduction}

\begin{itemize}
    \item The fairness problem in high-stakes ML decision systems: hiring,
          healthcare, lending.
    \item Positive-class scarcity as a structurally distinct regime: when
          historical inequality limits the supply of minority-positive training
          examples.
    \item Why this regime breaks conventional fairness interventions at the
          data level --- the gap motivating this work.
    \item The FORGE approach: coupling synthetic data generation to worst-group
          validation feedback via a REINFORCE agent in PCA space.
    \item Relationship to the conference paper (Pigou et al., IEEE CASSE 2025)
          and the journal paper (under review, Neurocomputing).
    \item Thesis contributions C1--C3 and chapter organisation.
\end{itemize}

% ===========================================================================
% Chapter 2
% ===========================================================================
\section{Background and Related Work}

\subsection{Group Fairness Metrics and Worst-Group Risk}
Equal opportunity, equalized odds, demographic parity, worst-group BCE loss.
Mathematical definitions and relationships. Why worst-group loss is preferred
as a training signal over threshold-sensitive fairness metrics.

\subsection{In-Processing Fairness Methods}
GroupDRO, FLB, adversarial debiasing, constraint-based methods. Why these
require modifying the classifier objective and what that costs in terms of
portability across downstream tasks.

\subsection{Data-Level Fairness Methods}
Instance reweighting, OT Repair, SMOTE and its lineage (ADASYN,
Borderline-SMOTE, cost-sensitive variants). Why interpolation-based
oversampling amplifies noise when the minority-positive pool is depleted.

\subsection{Generative Methods for Fairness}
CTGAN, TabFairGAN, FairFinGAN, FairTabDDPM, CuTS. Why distributional
training objectives without downstream classifier feedback cannot adapt to the
feature-space regions responsible for subgroup degradation.

\subsection{Reinforcement Learning for Data Augmentation and Valuation}
DVRL, Yang clinical fairness framework, EMRIL, and Pigou et al.\ (IEEE CASSE
2025). The conference paper established feasibility of RL-guided generation
for class-level imbalance; this thesis extends it to group fairness under
scarcity with a worst-group reward and PCA-compressed generation space.

\subsection{Summary and Gap}
No prior work trains an RL policy to generate synthetic minority-positive
samples with reward defined as worst-group validation loss on a downstream
classifier under the positive-class scarcity regime.

% ===========================================================================
% Chapter 3
% ===========================================================================
\section{The Positive-Class Scarcity Regime \contrib{C2}}

\subsection{Formal Definition}
Define DA+ (disadvantaged-group positive training examples). Characterise the
scarcity regime in terms of DA+ relative to training set size. Argue that at
DA+$\approx$43, the minority-positive pool is structurally too small to
support reliable reweighting or interpolation.

\subsection{Why Reweighting Methods Plateau}
\experiment{EXP-018 PCA-space baselines, census and Capture-24. GroupDRO and
SMOTE reduce EO but plateau at 0.078--0.114. Structural argument: with 43
minority-positive examples, reweighting amplifies noise rather than capturing
distributional structure.}

\subsection{Why Naive Generative Methods Fail}
\experiment{Spatial analysis of synthetic point placement in PCA space for
CTGAN, FairTabDDPM, SMOTE, and FORGE using existing run outputs and
\texttt{visualize\_generation.ipynb}. CTGAN (EO 0.328--0.400) and FairTabDDPM
(EO 0.151--0.313) place synthetic points without concentration near the
disadvantaged-positive cluster. Centroid drift analysis
(\texttt{plot\_centroid\_drift.py}) shows FORGE's synthetic centroid
progressively approaches the real disadvantaged-positive centroid while
untargeted methods do not. Failure mode: without downstream feedback,
generative methods cannot target the PCA subregion responsible for the
fairness gap.}

\subsection{Raw Feature Space Comparison}
\experiment{Table 4 --- FORGE (PCA-compressed) vs baselines trained on raw
features. Census: FORGE (EO 0.031) matches or beats best raw baselines (OT
Repair 0.033, FLB 0.036) while maintaining better utility than OT Repair (AUC
0.879 vs 0.863). Demonstrates that reward-guided generation in a compressed
space is not disadvantaged relative to uncompressed baselines.}

% ===========================================================================
% Chapter 4
% ===========================================================================
\section{FORGE Methodology \contrib{C1}}

\subsection{Framework Overview}
The three-model system: alpha classifier (real data baseline, trained once),
beta classifier (augmented, retrained from scratch each episode), REINFORCE
agent (generates synthetic minority-positive samples via delta actions in PCA
space).

\subsection{Environment and State Representation}
PCA-compressed feature space, delta-action walk, radius clipping, state
encoding (walk position + normalised progress). Justification for PCA:
retains dominant variance structure while making the generation search
tractable. Alternative spaces (raw features, VAE embeddings) discussed and
ruled out.

\subsection{Reinforcement Learning Agent}
Two-layer FFNN policy with independent Gaussian action heads per PCA
dimension, entropy bonus, Xavier initialisation. Why REINFORCE over PPO or
bandit alternatives: no differentiable environment model, trajectory-level
feedback structure.

\subsection{Reward Design}
Global sigmoid reward: $R = \sigma\!\left(k(\mathcal{L}^\text{worst}_\alpha -
\mathcal{L}^\text{worst}_\beta)\right)$. Why worst-group BCE rather than EO
directly (continuous, threshold-independent signal). Why uniform reward
distribution across trajectory steps (no per-step credit is available). Why
$\gamma = 1.0$. \\
\experiment{Ablation evidence: DVRL local reward (EXP-007/008) degrades
performance on both datasets; global-only reward confirmed as the reference
design.}

\subsection{Model Selection and Evaluation}
Best checkpoint retained by episode reward throughout training. Final
evaluation on held-out test set: EO, EOd, AUC, F1w.

% ===========================================================================
% Chapter 5
% ===========================================================================
\section{Experimental Setup}

\subsection{Datasets}
Census Income (natural positive-class scarcity, sex, DA+$=$43), Capture-24
(injected scarcity via group-specific subsampling, sex, DA+$=$45), ACS
Employment (disability framing, DA+$=$43 --- used as additional evidence in
Chapter 7, not a standalone contribution claim). DA+ log table and dataset
viability summary.

\subsection{Baselines}
Six baselines spanning reweighting (GroupDRO, OT Repair), in-processing (FLB),
interpolation oversampling (SMOTE), and generative synthesis (CTGAN,
FairTabDDPM). Two comparison settings: PCA-10 feature space (Table 3, primary)
and raw feature space (Table 4, supporting).

\subsection{Implementation Details}
FFNN architecture [32, 16], seeds, episode budgets. All hyperparameters used
in Chapters 6--7 are determined by the grid search in Chapter 6. Hardware and
software configuration.

% ===========================================================================
% Chapter 6
% ===========================================================================
\section{Hyperparameter Analysis \contrib{C3}}

\textit{Primary thesis contribution beyond the journal paper. The paper reports
the winning configuration; this chapter analyses the full landscape and
explains why particular configurations work.}

\subsection{Grid Search Design}
Four principal parameters: sigmoid sharpness $k \in \{0, 3, 5, 10\}$, PCA
dimensionality $d \in \{5, 10, 15\}$, classifier epochs per episode $\in
\{10, 20, 30\}$, synthetic data ratio $\in \{0.2, 0.4, 0.6\}$. Full
cartesian grid, 3 seeds per config. Census: 324 runs (EXP-021); Capture-24:
243 runs (EXP-025). Selection criterion: lowest mean $\beta$-EO, utility guard
$\beta$-AUC $\geq \alpha$-AUC $- 0.02$.

\subsection{Marginal Parameter Sensitivity}
\experiment{Sensitivity figure (four-panel) --- mean $\beta$-EO vs each
parameter value, marginalised over all other parameter combinations, one line
per dataset. Generated from EXP-021/025 grid data. Shows $k$ and classifier
epochs dominate; PCA dimensionality and synthetic ratio are secondary.}

\subsection{Interaction Effects}
\experiment{2D heatmaps from grid data --- $k \times \text{ep}$ and $k \times
\text{pca}$ combinations, marginalised over remaining parameters, per dataset.
Reveals whether high $k$ requires high epochs to be effective, and whether
optimal $k$ is stable across PCA choices.}

\subsection{Reward Signal Analysis and Cross-Dataset Parameter Differences}
\experiment{Episode return variance and deadzone fraction by $k$ and dataset,
computed from per-episode \texttt{metrics.csv} files across all grid runs.
Hypothesis: Capture-24's higher reward variance predicts $k{=}3$ (softer
sigmoid tolerates a noisy reward); census's lower variance supports $k{=}5$.
This provides a mechanistic explanation for why optimal $k$ differs across
datasets, and establishes episode return variance as a practical diagnostic for
selecting $k$ without a full grid search.}

\subsection{Stage 2: Secondary Parameter Random Search}
\experiment{EXP-022 --- 20 random configurations over secondary optimisation
parameters (classifier learning rate, policy learning rate, action scale
$\eta$, classifier optimiser, policy optimiser), conditioned on the best
principal configuration from Section 6.2. Random scatter figure ($\beta$-EO
vs AUC, 20 points). Selected values fill the implementation table.}

\subsection{Best Configurations and Summary}
Final selected configurations per dataset with rationale. These configurations
are used in all Chapter 7 comparisons.

% ===========================================================================
% Chapter 7
% ===========================================================================
\section{Empirical Evaluation \contrib{C1}}

\subsection{Policy Learning Verification}
\experiment{Training curves (\texttt{fig\_training\_curves\_nophase2.png},
existing) --- episode return over training, mean $\pm$ std across seeds for
both datasets. Centroid drift (\texttt{plot\_centroid\_drift.py}) --- L2
distance and cosine similarity between synthetic centroid and real
disadvantaged-positive centroid over training episodes, mean $\pm$ std.
Together verify: the policy learns a non-trivial generation strategy; a static
or random policy would produce a flat return profile and no directional
centroid drift.}

\subsection{Test-Set Generalisation}
\experiment{Generalisability curves (\texttt{check\_run.py --interval}) ---
test-set EO ($\downarrow$) and AUC ($\uparrow$) at snapshot intervals over
training episodes. Confirms that validation reward improvement transfers to
held-out test performance; no overfitting to the validation partition.}

\subsection{Comparative Performance --- PCA Feature Space}
\experiment{Table 3 (primary results) --- FORGE vs all six baselines, same
PCA-10 feature space, census and Capture-24. FORGE achieves lowest EO on
census (0.031 $\pm$ 0.018 vs FLB 0.039) while maintaining near-baseline
utility (AUC 0.879). Capture-24 provisional best pending full grid completion.}

\subsection{Comparative Performance --- Raw Feature Space}
\experiment{Table 4 (supporting) --- FORGE in PCA-compressed space vs
baselines trained on raw features. Census: FORGE (0.031) achieves lowest or
tied-lowest EO against baselines with full feature access, while maintaining
better AUC than OT Repair. Capture-24: $\alpha$-EO discrepancy between feature
spaces noted; results interpreted via relative EO reduction rather than
absolute values.}

\subsection{ACS Employment --- Additional Evidence}
\experiment{EXP-027 two-config run (best census config $k{=}5$, pca$=10$,
ep$=30$; best Capture-24 config $k{=}3$, pca$=10$, ep$=10$) on ACS
Employment, disability framing, DA+$=$43. Demonstrates C1 holds on a third
dataset with a non-demographic protected attribute. Not a generalisation claim;
presented as supporting evidence that the method produces meaningful EO
reduction beyond the two primary datasets.}

% ===========================================================================
% Chapter 8
% ===========================================================================
\section{Discussion}

\subsection{What Governs FORGE's Effectiveness \contrib{C3}}
Synthesise Chapter 6 findings. Reward sharpness $k$ is the dominant parameter;
its optimal value tracks with episode return variance, a dataset-level property
measurable during training without a full grid search. Classifier epochs
determine reward informativeness: too few leaves beta underconverged; too many
overfits beta to the current synthetic batch. Practitioners can use deadzone
fraction and return variance as early diagnostics for configuration selection.

\subsection{Seed Variance Analysis}
Census ($\beta$-EO std $\pm$0.018) vs Capture-24 ($\pm$0.051). Analyse
per-seed $\alpha$-EO and deadzone fraction to determine whether variance
originates in the data split or the RL optimisation. Connects to the scarcity
regime characterisation in Chapter 3.

\subsection{Failure Cases}
Which seeds fail to improve $\beta$-EO over $\alpha$-EO? Is there a
leading indicator --- deadzone fraction, episode return trajectory --- that
predicts failure early in training?

\subsection{Relationship to Prior Work}
Explicit comparison to Pigou et al.\ (IEEE CASSE 2025): same core idea
(RL-guided generation), different problem (group fairness under scarcity vs
class-level imbalance), different reward (worst-group BCE vs decision-boundary
proximity), different feature space (PCA-compressed), formal ablation studies
and multi-dataset evaluation added.

\subsection{Limitations}
Trajectory-level credit assignment (all steps share identical reward; no
per-step attribution). Single classifier architecture (FFNN); method is
data-level but untested with other downstream models. Binary protected
attributes only. Single scarcity level evaluated. Computational cost of
per-episode classifier retraining ($\approx$50 minutes per seed per dataset).

\subsection{Future Work}
Evolutionary search (CMA-ES or NES) to address trajectory-level credit
assignment. Multi-group and intersectional fairness extensions. Learned latent
spaces (VAE, diffusion) as alternatives to PCA. Warm-starting beta classifier
to reduce per-episode training cost.

% ===========================================================================
% Chapter 9
% ===========================================================================
\section{Conclusion}

\begin{itemize}
    \item Restate contributions C1--C3 and the evidence supporting each.
    \item The central finding: reward-guided generation is a viable fairness
          intervention under positive-class scarcity where both reweighting
          and naive generative methods demonstrably fail.
    \item Practical implication of C3: episode return variance and deadzone
          fraction are actionable diagnostics for reward sharpness selection,
          reducing the need for exhaustive grid search in new applications.
    \item Directions for future work.
\end{itemize}

\vspace{1em}
\hrule
\vspace{1em}

% ===========================================================================
% Contribution--Experiment Mapping
% ===========================================================================
\section*{Contribution--Experiment Mapping}

\begin{table}[h]
\centering
\renewcommand{\arraystretch}{1.3}
\begin{tabular}{p{0.08\linewidth} p{0.38\linewidth} p{0.38\linewidth} p{0.08\linewidth}}
\toprule
\textbf{Contrib.} & \textbf{Key experiments} & \textbf{Data source} & \textbf{Chapter} \\
\midrule
C1 & Main results tables (PCA + raw), training curves, centroid drift, generalisation curves, DVRL ablation & EXP-001/002, EXP-007/008, EXP-018, raw baselines, plot scripts & 4, 7 \\
\midrule
C2 & Spatial failure mode analysis (PCA visualisation), baseline plateau analysis, raw feature comparison & EXP-018, existing run outputs, \texttt{visualize\_generation.ipynb} & 3 \\
\midrule
C3 & Full grid landscape, sensitivity figure, interaction heatmaps, reward signal analysis, random search & EXP-021 (census, 324 runs), EXP-025 (Capture-24, 243 runs), EXP-022 & 6, 8.1 \\
\bottomrule
\end{tabular}
\end{table}

\end{document}
